\documentclass[12pt,a4paper,openany]{book}


\usepackage{lmodern}
\usepackage{xcolor}
\usepackage{xcolor}
\definecolor{vert1}{rgb}{0.0,0.3.9,0.0}
\definecolor{bleu}{rgb}{0,0,0.5}
\definecolor{bleu3}{rgb}{1,0.2,0.2}
\definecolor{grisgris}{gray}{0.4}
\definecolor{grisclair}{HTML}{E7E7E7}
\definecolor{grisfonce}{HTML}{A5A5A5}
\definecolor{rougeUPS}{rgb}{0.6, 0.3, 0.3}

\fboxsep =0pt \parindent =0pt\parskip =12pt



\usepackage[utf8]{inputenc}
\usepackage[T1]{fontenc}
\usepackage[francais]{babel}
\usepackage[top=1.7cm, bottom=1.7cm, left=1.7cm, right=1.7cm]{geometry}
\usepackage{verbatim}
\usepackage[urlbordercolor={1 1 1}, linkbordercolor={1 1 1}, linkcolor=vert1, urlcolor=bleu, colorlinks=true]{hyperref}
\usepackage{tikz} %Vectoriel
\usepackage{listings}
\usepackage{fancyhdr}
\usepackage{multido}
\usepackage{amssymb}

\newcommand{\titre}{Architecture des systèmes Informatiques --- TD}

\newcommand{\pole}{}
\newcommand{\sigle}{archi}

\newcommand{\semestre}{3}

\definecolor{gris1}{gray}{0.40}
\definecolor{gris2}{gray}{0.55}
\definecolor{gris3}{gray}{0.65}
\definecolor{gris4}{gray}{0.50}
\definecolor{vert}{rgb}{0,0.4,0}
\definecolor{violet}{rgb}{0.65, 0.2, 0.65}
\definecolor{bleu1}{rgb}{0,0,0.8}
\definecolor{bleu2}{rgb}{0,0.2,0.6}
\definecolor{bleu3}{rgb}{0,0.2,0.2}
\definecolor{rouge}{HTML}{F93928}


\lstdefinelanguage{algo}{%
   morekeywords={%
    %%% couleur 1
		importer, programme, glossaire, fonction, procedure, constante, type, 
	%%% IMPORT & Co.
		si, sinon, alors, fin, tantque, debut, faire, lorsque, fin lorsque, 
		declenche, declencher, enregistrement, tableau, retourne, retourner, =, pour, a,
		/=, <, >, traite,exception, 
	%%% types 
		Entier, Reel, Booleen, Caractere, Réél, Booléen, Caractère,
	%%% types 
		entree, maj, sortie,entrée,
	%%% types 
		et, ou, non,
	},
  sensitive=true,
  morecomment=[l]{--},
  morestring=[b]',
}

\lstset{language=algo,
    %%% BOUCLE, TEST & Co.
      emph={importer, programme, glossaire, fonction, procedure, constante, type},
      emphstyle=\color{bleu2},
    %%% IMPORT & Co.  
	emph={[2]
		si, sinon, alors, fin , tantque, debut, faire, lorsque, fin lorsque, 
		declencher, retourner, et, ou, non,enregistrement, retourner, retourne, 
		tableau, /=, <, =, >, traite,exception, pour, a
	},
      emphstyle=[2]\color{bleu1},
    %%% FONCTIONS NUMERIQUES
      emph={[3]Entier, Reel, Booleen, Caractere, Booléen, Réél, Caractère},
      emphstyle=[3]\color{gris1},
    %%% FONCTIONS NUMERIQUES
      emph={[4]entree, maj, sortie, entrée},	
      emphstyle=[4]\color{gris1},
}
\lstdefinelanguage{wl}{%
   morekeywords={%
    %%% couleur 1
		importer, programme, glossaire, fonction, procedure, constante, type, 
	%%% IMPORT & Co.
		si, sinon, alors, fin, TANTQUE, tantque, FIN, PROCEDURE, debut, faire, lorsque, 
		fin lorsque, declenche, declencher, enregistrement, tableau, retourne, retourner, =, 
		/=, <, >, traite,exception, 
	%%% types 
		Entier, Reel, Booleen, Caractere, Réél, Booléen, Caractère,
	%%% types 
		entree, maj, sortie,entrée,
	%%% types 
		et, ou, non,
	},
  sensitive=true,
  morecomment=[l]{//},
  morestring=[b]',
}

\lstset{language=wl,
    %%% BOUCLE, TEST & Co.
      emph={importer, programme, glossaire, fonction, procedure, constante, type},
      emphstyle=\color{bleu2},
    %%% IMPORT & Co.  
	emph={[2]
		si, sinon, alors, fin , tantque, debut, faire, lorsque, fin lorsque, 
		declencher, retourner, et, ou, non,enregistrement, retourner, retourne, 
		tableau, /=, <, =, >, traite,exception
	},
      emphstyle=[2]\color{bleu1},
    %%% FONCTIONS NUMERIQUES
      emph={[3]Entier, Reel, Booleen, Caractere, Booléen, Réél, Caractère},
      emphstyle=[3]\color{gris1},
    %%% FONCTIONS NUMERIQUES
      emph={[4]entree, maj, sortie, entrée},	
      emphstyle=[4]\color{gris1},
}
\lstdefinelanguage{css}{%
   morekeywords={%
    %%% couleur 1
		background, image, repeat, position, index, color, border, font, 
		size, url, family, style, variant, weight, letter, spacing, line, 
		height, text, decoration, align, indent, transform, shadow, 
		background, image, repeat, position, index, color, border, font, 
		size, url, family, style, variant, weight, letter, spacing, line, 
		height, text, decoration, align, indent, transform, shadow, 
		vertical, align, white, space, word, spacing,attachment, width, 
		max, min, margin, padding, clip, direction, display, overflow,
		visibility, clear, float, top, right, bottom, left, list, type, 
		collapse, side, empty, cells, table, layout, cursor, marks, page, break,
		before, after, inside, orphans, windows, azimuth, after, before, cue, 
		elevation, pause, play, during, pitch, range, richness, spek, header, 
		numeral, punctuation, rate, stress, voice, volume,
	%%% types 
		left, right, bottom, top, none, center, solid, black, blue, red, green,
	},
  sensitive=true,
  sensitive=true,
  morecomment=[s]{/*}{*/},
  morestring=[b]',
}
\lstset{language=css,
    %%% BOUCLE, TEST & Co.
      emph={
		background, image, repeat, position, index, color, border, font, 
		size, url, family, style, variant, weight, letter, spacing, line, 
		height, text, decoration, align, indent, transform, shadow, 
		background, image, repeat, position, index, color, border, font, 
		size, url, family, style, variant, weight, letter, spacing, line, 
		height, text, decoration, align, indent, transform, shadow, 
		vertical, align, white, space, word, spacing,attachment, width, 
		max, min, margin, padding, clip, direction, display, overflow,
		visibility, clear, float, top, right, bottom, left, list, type, 
		collapse, side, empty, cells, table, layout, cursor, marks, page, break,
		before, after, inside, orphans, windows, azimuth, after, before, cue, 
		elevation, pause, play, during, pitch, range, richness, spek, header, 
		numeral, punctuation, rate, stress, voice, volume,
	  },
      emphstyle=\color{bleu2},
    %%% FONCTIONS NUMERIQUES
      emph={[3]
		left, right, bottom, top,none, solid, black, blue, green,
		  },
      emphstyle=[3]\color{bleu3},
    %%% FONCTIONS NUMERIQUES
}

\lstset{language=SQL,
    %%% BOUCLE, TEST & Co.
      emph={INSERT, UPDATE, DELETE, WHERE, SET, GROUP, BY, ORDER, REFERENCES},
      emphstyle=\color{bleu2},
    %%% IMPORT & Co.  
	emph={[2]
		if, end, begin, then, for, each, else, after, of, on, to
	},
      emphstyle=[2]\color{bleu1},
    %%% FONCTIONS NUMERIQUES
      emph={[3]Entier, Reel, Booleen, Caractere, Booléen, Réél, Caractère},
      emphstyle=[3]\color{gris1},
    %%% FONCTIONS NUMERIQUES
      emph={[4]entree, maj, sortie, entrée},	
      emphstyle=[4]\color{gris1},
}
\lstdefinelanguage{ARM}{%
   morekeywords={%
   ADD, SUB, MOV, MUL, RSB,CMP, BLS, BLE, B,BHI,LDR,
   BGE, RSBLT, BGT, BEQ, BNE,BLT,BHS,STR,STRB
	},
  sensitive=true,
  morecomment=[l]{@},
  morestring=[b]',
}

\lstset{ % general style for listings 
   numbers=left 
   , literate={é}{{\'e}}1 {è}{{\`e}}1 {à}{{\`a}}1 {ê}{{\^e}}1 {É}{{\'E}}1 {ô}{{\^o}}1 {€}{{\euro}}1{°}{{$^{\circ}$}}1 {ç}{ {c}}1 {ù}{u}1
	, extendedchars=\true
   , tabsize=2 
   , frame=l
   , framerule=1.1pt
   , linewidth=520px
   , breaklines=true 
   , basicstyle=\footnotesize\ttfamily 
   , numberstyle=\tiny\ttfamily 
   , framexleftmargin=0mm 
   , xleftmargin=0mm 
   , captionpos=b 
	, keywordstyle=\color{bleu2}
	, commentstyle=\color{vert}
	, stringstyle=\color{rouge}
	, showstringspaces=false
	, extendedchars=true
	, mathescape=true
} 
%	\lstlistoflistings
%	\addcontentsline{toc}{part}{List of code examples}
 %prise en charge du langage algo
\date{\today}

\makeindex
\lfoot{Université Toulouse III -- Paul Sabatier}
\rfoot{\sigle\semestre}
%\rfoot{}
\cfoot{}
\makeglossary
\makeatletter
\def\clap#1{\hbox to 0pt{\hss #1\hss}}%
\def\ligne#1{%
\hbox to \hsize{%
\vbox{\centering #1}}}%
\def\haut#1#2#3{%
\hbox to \hsize{%
\rlap{\vtop{\raggedright #1}}%
\hss
\clap{\vtop{\centering #2}}%
\hss
\llap{\vtop{\raggedleft #3}}}}%
\def\bas#1#2#3{%
\hbox to \hsize{%
\rlap{\vbox{\raggedright #1}}%
\hss \clap{\vbox{\centering #2}}%
\hss
\llap{\vbox{\raggedleft #3}}}}%
\def\maketitle{%
\thispagestyle{empty}\vbox to \vsize{%
\haut{}{\@blurb}{}

\vfill
\vspace{1cm}
\begin{flushleft}
\usefont{OT1}{ptm}{m}{n}
\huge \@title
\end{flushleft}
\par
\hrule height 4pt
\par
\begin{flushright}
\usefont{OT1}{phv}{m}{n}
\Large \@author
\par
\end{flushright}
\vspace{1cm}
\vfill
\vfill
\bas{}{\@location, le \@date}{}
}%
\cleardoublepage
}
\def\date#1{\def\@date{#1}}
\def\author#1{\def\@author{#1}}
\def\title#1{\def\@title{#1}}
\def\location#1{\def\@location{#1}}
\def\blurb#1{\def\@blurb{#1}}
\date{\today}
\author{}
\title{}
\location{Amiens}\blurb{}
\makeatother
\title{\titre}
\author{Semestre \semestre}

\location{Toulouse}
\blurb{%
Université Toulouse III -- Paul sabatier\\
L2 Informatique\\
}%



%\title{Cours \\ \titre}
%\date{\today\\ Semestre \semestre}

%\lhead{Cours: \titre}
%\chead{}
%\rhead{\thepage}

%\lfoot{Université Paul Sabatier Toulouse III}
%\cfoot{\thepage}
%\rfoot{\sigle\semestre}

\pagestyle{fancy}
\renewcommand{\chaptermark}[1]{\markboth{\bsc{\chaptername~\thechapter{} :} #1}{}}
\renewcommand{\sectionmark}[1]{\markright{\thesection{ #1}}}
\renewcommand{\headrulewidth}{0.3pt}
\renewcommand{\footrulewidth}{0.3pt}

\fancyhf{}
\fancyhead[LE]{\leftmark}
\fancyhead[RO]{\rightmark}
\fancyfoot[LE,RO]{--~\thepage~--}
\fancyfoot[LO]{\titre{}~~---~~\sigle{}\semestre{}}
\fancyfoot[RE]{Antoine de \bsc{Roquemaurel}}

%% Cas des premières pages de chapitre
\fancypagestyle{plain}{%
	\fancyhf{}%
	\fancyfoot[L]{\titre{}~~---~~\sigle{}\semestre{}}
	\fancyfoot[R]{--~\thepage~--}
	\renewcommand{\headrulewidth}{0pt}
	\renewcommand{\footrulewidth}{0.3pt}
}
\makeatletter
\renewcommand*{\lstlistlistingname}{Liste des codes sources}
\renewcommand\listoffigures{%
    \chapter{\listfigurename}%
      \@mkboth{\MakeUppercase\listfigurename}%
              {\MakeUppercase\listfigurename}%
       \@starttoc{lof}%
    }
    \renewcommand\listoftables{%
    \chapter{\listtablename}%
    \@mkboth{\MakeUppercase{\listtablename}}%
            {\MakeUppercase{\listtablename}}%
    \@starttoc{lot}
    }

    \renewcommand\lstlistoflistings{%
    \begingroup
    \chapter{\lstlistlistingname}%
    \parskip\z@\parindent\z@\parfillskip \z@ \@plus 1fil%
    \@starttoc{lol}%
    \endgroup
    }
	\makeatother

\newcommand{\remarque}[1]{
	\begin{center}
	\medskip
	\colorbox{remarque}{
		\begin{minipage}{0.85\textwidth}\medskip\includegraphics[height=10px]{images/remarque.png} #1 \medskip\end{minipage}
	}
	\medskip
	\end{center}
}

\newcounter{exemples}

\newenvironment{exemple}[1]{
   \vspace{-2mm}

\refstepcounter{exemples}
   \begin{center}
	\medskip
      \begin{minipage}{0.9\linewidth}
}{%
~
      \end{minipage}
   \end{center}~
   \vspace{-2mm}
}%

\newcommand{\captionExemple}[1]{
	\begin{center}{\bsc{Exemple} \thechapter.\arabic{exemples}~--~}#1\end{center}
}

\DeclareTextFontCommand{\policeGlossaire}{\fontfamily{lmss}\selectfont}
\DeclareTextFontCommand{\policePackage}{\fontfamily{phv}\selectfont}
\DeclareTextFontCommand{\policeTitre}{\fontfamily{ptm}\selectfont}
\newcommand{\policeCode}[1]{\texttt{#1}}

\newcommand{\sectionfont}{%
	\fontencoding{\encodingdefault}%
	\fontfamily{pag}%
	\fontseries{bc}%
	\fontshape{n}%
	\selectfont
}

% numéro du chapitre
\DeclareFixedFont{\chapnumfont}{T1}{phv}{b}{n}{80pt}
% pour le mot « Chapitre »
\DeclareFixedFont{\chapchapfont}{T1}{phv}{b}{n}{16pt}
% pour le titre
\DeclareFixedFont{\chaptitfont}{T1}{phv}{b}{n}{24.88pt}


\makeatletter
\def\thickhrulefill{\leavevmode \leaders \hrule height 1ex \hfill \kern \z@}
%% \chapter
\def\@makechapterhead#1{%
  \reset@font
  \parindent \z@
  \vspace*{10\p@}%
  \hbox{%
    \vbox{%
      \advance\hsize by -2cm
      \hrule height 0.4pt depth 0pt width \hsize
      \par
      \vskip 6pt%
      \hspace{20pt}%
      \parbox{420pt}{%
        \LARGE \bfseries #1
		}%
      \par
      \vskip 6pt%
      \hspace{20pt}%
      \hrule height 0.4pt depth 0pt width \hsize
	  \vspace{-30pt}
      }%
    \vbox{%
      \hsize=1.5cm%
      \begin{tabular}{c}
        \scshape \large \strut \@chapapp{} \\
        \colorbox{black}{\vbox{\hbox{\vbox to 1mm{}}\hbox{
			\color{white} \LARGE \bfseries \hspace{1mm}\thechapter\hspace{1mm}
		}\hbox{\vbox to 2cm{}}}}%
      \end{tabular}%
      }%
    }%
  \vskip 20\p@
}
%% \chapter*
\def\@makeschapterhead#1{%
  \reset@font
  \parindent \z@
  \vspace*{10\p@}%
  \hbox{%
    \vbox{%
      \advance\hsize by -0cm
      \hrule height 0.4pt depth 0pt width \hsize
      \par
      \vskip 6pt%
      \hspace{20pt}%
      \parbox{420pt}{%
        \LARGE \bfseries #1
		}%
      \par
      \vskip 6pt%
      \hspace{20pt}%
      \hrule height 0.4pt depth 0pt width \hsize
      }%
    }%
  \vskip 20\p@

}

\newlength{\sectiontitleindent}
\newlength{\subsectiontitleindent}
\newlength{\subsubsectiontitleindent}
\setlength{\sectiontitleindent}{-1cm}
\setlength{\subsectiontitleindent}{-.5cm}
\setlength{\subsubsectiontitleindent}{-.25cm}

\renewcommand{\section}{%
	\@startsection%
	{section}%
	{1}%
	{\sectiontitleindent}%
	{-3.5ex plus -1ex minus -.2ex}%
	{2.3ex plus.2ex}%
	{\sectionfont\Large}
}
\renewcommand{\subsection}{%
	\@startsection%
	{subsection}%
	{2}%
	{\subsectiontitleindent}%
	{-3.5ex plus -1ex minus -.2ex}%
	{2.3ex plus.2ex}%
	{\sectionfont\large}
}

\renewcommand{\subsubsection}{%
	\@startsection%
	{subsubsection}%
	{3}%
	{\subsubsectiontitleindent}%
	{-3.5ex plus -1ex minus -.2ex}%
	{2.3ex plus.2ex}%
	{\sectionfont\normalsize}
}

\makeatother




\begin{document}
	\setcounter{tocdepth}{2}
	\setcounter{secnumdepth}{3}
	\maketitle
	\tableofcontents
	\chapter{Numérotation et codage\\ TD optionnel}
		\section{Numérotation}
		\subsection{Réaliser l'opération suivante en binaire : $(1101011-10110)\times 11001$}
		\begin{eqnarray*}
			(110\;1011)_2 &=& (107)_{10}\\
			(1\;0110)_2&=&(22)_{10}\\
			(1\;1001)_2&=&(22)_{10}
		\end{eqnarray*}
		\begin{center}
		\begin{tabular}{cccccccc}
			&1&1&0&1&0&1&1\\
			$-$&$$&&1&0&1&1&0\\
			\hline 
			&1&0&1&0&1&0&1\\
		\end{tabular} Ou 
		\begin{tabular}{cccccccc}
			&1&1&0&1&0&1&1\\
			+&1&1&1&0&1&0&1\\
			\hline
			&1&0&1&0&1&0&1
		\end{tabular}
		\begin{eqnarray*}
			(1010101)_2&=&(85)_{10}
		\end{eqnarray*}
		\begin{tabular}{cccccccccccccc}
			&&&&&&1&0&1&0&1&0&1&\\
			&&&&&$\times$&&1&1&0&0&1&\\
			\hline
			&&&&&&1&0&1&0&1&0&1&\\
			&&&1&0&1&0&1&0&1&.&.&.&\\
			&&1&0&1&0&1&0&1&.&.&.&.&\\
			\hline
			1&0&0&0&0&0&1&0&0&1&1&0&1&
		\end{tabular}
	\end{center}
		\subsection{Réaliser les opérations suivantes en hexadécimal : $(389A+7293)-EB2$}
		\begin{center}
			
		\begin{tabular}{ccccc}
			&3&8&9&A\\
			+&7&2&9&3\\
			\hline
			&A&B&2&B\\
			$-$&E&B&2&B\\
			\hline
			&9&C&7&B
		\end{tabular}
		\begin{eqnarray*}
			(389A)_{16} &=& (0011\; 1000\; 1001\; 1010)_2\\
			(7293)_{16} &=& (0111\; 0010\; 1001\; 0011)_2\\
			(AB2B)_{16} &=& (1010\;1011\;0010\;1101)_2\\
			(EB2B)_{16} &=& (0000\;1110\;1011\;0010)_2\\
			(9C7B)_{16} &=& (1001\;1100\;0111\;1011)_2
		\end{eqnarray*}
		\end{center}

		\subsection{Effectuer les conversions ci-dessous}
			\subsubsection{$(1447.140625)_{10} = (??)_2 = (??)_{16}$}
			\begin{eqnarray*}
				1447\div 16 &=& 90\; R=7\\
				90 \div 16 &=&~5\; R=A\\
				5 \div 16 &=&  ~0 \; R=5\\ 
				(1447)_{10} &=& (5A7)_{16}\\\\
				0.140625 \times 16 &=& 2.25\\
				0.25 \times 16 &=& 4.00\\
				(0.140625)_{10} &=& (0.24)_{16}\\\\
				(1447.140625)_{10} = (5A7.24)_{16} &=& (0101\;1010\;0111.0010\;0100)_2
			\end{eqnarray*}
			\subsubsection{$(1111100101.01011)_{2} = (??)_{10} = (??)_{16}$}
			\begin{eqnarray*}
				(11\;1110\;0101.0\;1011)_2 &=& (3E5;58)_{16}\\
				3E5 &=& (3\times 16^2 + 14*16 + 5 + 5 \times 16^{-1} + 8 \times 16 ^{-2}\\
				&=& (997,34375)_{10}
			\end{eqnarray*}
		\section{Codage}
		\subsection{Codage de nombres entiers relatifs}
		On veut coder des entiers relatifs sur 16 chifffres binaires (deux octets).

		\begin{tabular}{| p{1.8cm}|p{4cm}|p{2cm}|p{3.7cm}|p{3.7cm}|}
			\hline
			& \multicolumn{2}{ c|}{\textbf{Valeur absolue}}  & \multicolumn{2}{ c |}{\textbf{Valeur relative}}\\
			\hline
			\textbf{Base 10} & \textbf{Base 2} & \textbf{Base 16} & \textbf{Valeur absolue + signe} & \textbf{Complément à 2}\\
			\hline
			$\ \ 35671$&$10000\;010\;0101\;0111$& 75B8 & Hors intervalle & Hors intervalle\\
			\hline
			$-32768$ & $1000\; 0000\; 0000\; 0000$ & $8000$ & Hors intervalle & $1000\; 0000\; 0000$\\
			\hline
			$\ \ 46443$& $1011\; 0101\; 0110\; 1011$ & B56B & Hors intervalle & Hors intervalle\\
			\hline
			$-19536$& $0100\; 1100\; 0110\;0100$ & 4C64 & $1100\; 1100\; 0110\; 0100$ & $1011\; 0011\; 1001\; 1100$\\
			\hline
			$-19040$ & $0100\; 1010\; 0110\; 0000$ & 4A60 & $1100\; 1010\; 0110\; 0000$ & $1011\; 0101\; 1010\; 0000$\\
			\hline
		\end{tabular}
		\subsubsection{Calcul de la valeur absolue en base 2}
		\begin{eqnarray*}
			3567 \div 16 &=& 229 \textrm{ et reste 7}\\
			2229 \div 16 &=& 139 \textrm{ et reste 5}\\
			139 \div 16 &=& 8 \textrm{ et reste B}\\
			8 \div 16 &=&0 \textrm{ et reste 8}\\
		\end{eqnarray*}
		$$\Rightarrow (35671)_{10} = (8B57)_{16} = (1000\; 1010\; 0101\; 0111)_2$$
		\remarque{
		Afin de pouvoir représenter un nombre celui-ci ne doit pas dépasser un certain intervalle : 

\begin{description}
	\item[Entier naturel] $[0;2^{16-1}] = [0;65535]$
	\item[Valeur absolue + signe] $[-2^{15-1};2^{15-1}] = [-16384 ; 16384]$
	\item[Complément à deux] $[-2^{15};2^{15}] = [-32768 ; 32768]$
\end{description}
		}
		\subsection{Convertir un nombre flottant en décimal}
		\begin{center}
			\begin{tabular}{|c|c|c|}
				\hline
				$0$ & $1000\; 0001$ & $1110\; 0000\; 0000\; 0000\; 0000\; 0000$\\
				\hline
			\end{tabular}
		\end{center}
	$$C = E + biais$$
	En simple précision biais = 127.

	\begin{eqnarray*}
		S &=& 0 \rightarrow \textrm{ positif}\\
		C &=& 129\\
		E &=& C - 127 = 2\\
		1.M &=& 1.111 \Rightarrow 1.111 \times 2^2 = 111.1 \times 2^0 = (7.5)_{10}\\
	\end{eqnarray*}
	\newpage	
	\subsection{Convertir un nombre décimal en flottant}
	\begin{eqnarray*}
		(35.5)_{10} &=& ? \\
		100011.1 &=& 10000111 \times 2^5\\
		\textrm{Nombre positif donc } S &=& 0\\
		E &=& 5\\
		C &=& E + 127 = 132 = 128+4 = (10000100)_2\\
		1.M &=& 1.00011\\
		M &=& 0011\\
	\end{eqnarray*}
	
		\begin{center}
			\begin{tabular}{|c|c|c|}
				\hline
				$0$ & $1000\; 0001$ & $0100\; 0011\; 0000\; 0000\; 0000\; 0000$\\
				\hline
			\end{tabular}
		\end{center}
		

		\chapter{Algèbre de Bool}
	\section{Table de vérité des opérateurs classiques}
	\subsection{Exercice 1}
		\subsubsection{A -- Démontrer que les opérateurs NAND et NOR sont des opérateurs complets}
		\begin{eqnarray*}
			\overline A &=& A|A\\
			A.B &=& \overline{\overline{A.B}} = \overline{\overline A + \overline B} = \overline{A/B} = (A | B) | (A | B)\\
			A+B &=& \overline{\overline{A+B}} = \overline{\overline A + \overline B} = (A|B) |(B |B)
		\end{eqnarray*}
		\begin{eqnarray*}
			\overline A = \overline{A+A} = A \downarrow A\\
			A . B = \overline{\overline{A.B}} = \overline{\overline A+ \overline B} = A \downarrow B \downarrow (\downarrow B \downarrow B)
		\end{eqnarray*}
		\subsubsection{B --}
		\begin{enumerate}
			\item $f(A,B,C,D)=\overline A B \overline D + B\overline C + A\overline C D= \overline{\overline{\overline A B \overline D + B \overline C + A \overline C D}} 
				= \overline{(\overline{\overline B\overline D).(\overline{A\overline C D)}}} =\\ ( (A|A)|B( (D|D))) | (B|(C|C)) | (A|(C|C)|D)$
			\item $f(A,B,C,D)=(A+B)(\overline C + \overline D)(\overline A + \overline B + \overline C) = \overline{\overline{(A+B)(\overline C + \overline B)(\overline A + \overline B + \overline C)}}=\\
				\overline{\overline A \overline B) + CD + ABC} = \overline{\overline A \overline B} . \overline{CD} . \overline{ABC} =\\
				((A|A)(B|B)|(C|D)|(A|B|C))|((A|A)(B|B)|(C|D)|(A|B|C)) =\\
				((A\downarrow B)\downarrow (C \downarrow C) \downarrow (D \downarrow D) \downarrow (A \downarrow A)\downarrow (B\downarrow B) \downarrow (C \downarrow C))$
		\end{enumerate}
	\subsection{Exercice 3}
	\subsubsection{1}
	$f(w,x,y,z) = \sum n(1,3,4,7,9,11,12,15) = yz+\overline x z + x\overline y \overline z$
	\\
	\begin{tabular}{c|c|c|c|c|}
		& 00&01&11&10\\
		\hline
		00&&1&1&\\
		\hline
		01&1&&1&\\
		\hline
		11&1&&1&\\
		\hline
		10&&1&1&\\
		\hline
	\end{tabular}
	\subsubsection{2}
	$f(w,x,y,z) = \sum m(0,1,3,6,9,13,15) = \overline w \overline x \overline y + \overline w \overline x z + wxz + w\overline y z+w\overline x y \overline z$\\
	\begin{tabular}{c|c|c|c|c|}
		& 00&01&11&10\\
		\hline
		00&1&1&1&\\
		\hline
		01&&&&1\\
		\hline
		11&&1&1&\\
		\hline
		10&&1&&\\
		\hline
	\end{tabular}
	\subsubsection{3}
	$f(w,x,y,z) = \sum m (0,1,5,7,8,10,14,15) = \overline w \overline x \overline y +  \overline w x z + wxy + w\overline x \overline  z$\\
	\begin{tabular}{c|c|c|c|c|}
		& 00&01&11&10\\
		\hline
		00&1&1&&\\
		\hline
		01&&1&1&\\
		\hline
		11&&&1&1\\
		\hline
		10&1&&&1\\
		\hline
	\end{tabular}

	\subsection{4}
	\begin{tabular}{c|c|c|c|c|}
		& 00&01&11&10\\
		\hline
	\end{tabular}
	\subsection{5}
	\begin{tabular}{c|c|c|c|c|}
		& 00&01&11&10\\
		\hline
	\end{tabular}
	\subsection{9}
	$f(w,x,y,z) = \Pi M(1,3,4,9,11,14) CI(w,x,y,z) = (x+\overline z)(\overline x + z)$\\
	\begin{tabular}{c|c|c|c|c|}
		& 00&01&11&10\\
		\hline
		00&&0&0&\\
		\hline
		01&0&&&*\\
		\hline
		11&*&*&&0\\
		\hline
		10&&0&0&\\
		\hline
	\end{tabular}
	\subsection{11}
	$f(w,x,y,z) = \sum m(0,1,3,5,6,7,11,13,14,15) CI(w,x,y,z) = 4\\
	f(w,x,y,z) = \overline y \overline w + yz + xy +xz$\\
	\begin{tabular}{c|c|c|c|c|}
		& 00&01&11&10\\
		\hline
		00&1&1&1&\\
		\hline
		01&*&1&1&1\\
		\hline
		11&&1&1&1\\
		\hline
		10&&&1&\\
		\hline
	\end{tabular}

	\chapter{Fonctions logiques}
	\section{Exercice 1 : Simplifications algébriques}
	\subsection{}
	\begin{tabular}{ccc||c}
		x&y&z & $F_1$\\
		0&0&0&	0\\
		0&0&1& 	1\\
		0&1&0& 	1\\
		0&1&1& 	1\\
		1&0&0& 	1\\
		1&0&1& 	1\\
		1&1&0& 	1\\
		1&1&1& 	1\\
	\end{tabular}
	\begin{eqnarray*}
		F_1 &=& (x+y+z).(\overline x + \overline y + z + xy + \overline x \overline y)\\
		&=& (x+y+z).(\overline{\overline x} + \overline{\overline y} + z + xy + \overline{\overline x \overline y})\\
		&=& (x+y+z)(\overline x + \overline y (1 + \overline x) + z +xy)\\
		&=& (x+y+z)(\overline x + \overline y + z + xy)\\
		&=& (x+y+z)(\overline{xy} + xy + z)\\
		&=& (x+y+z)(1+z)\\
		F_1 &=& x+y+z
	\end{eqnarray*}
	\subsection{}
	\begin{eqnarray*}
		F_2 &=& \sum m(0,4,6,7,14,15)\\
		&=& \overline x \overline y \overline z \overline w + \overline x y \overline z \overline w + \overline x y z \overline w + \overline x yzw + xyz\overline w + xyzw\\
		&=& \overline x \overline z \overline w(y+\overline y) +\overline x yz (w+\overline w) + xyz(w + \overline w)\\
		&=& \overline x \overline z \overline w + \overline x y z + xyz\\
		&=& \overline x \overline z \overline w + xy(x + \overline x)\\
		F_2 &=& \overline x \overline z \overline w + xy
	\end{eqnarray*}
	\section{Exercice 2 : Formes canoniques}
	\subsection{}
	$G_1 = I_1 + I_2$ avec $I_1 = \sum m (0,4,6)$ et $I_2 = \Pi M(1,4,5)$
	\begin{eqnarray*}
		G_1 &=& \sum m(0,4,6) + \Pi M(1,4,5)\\
		&=& \sum m(0,4,6) + \sum m (0,2,3,6,7)\\
		G_1 &=& \sum m (0,2,3,4,6,7) \Rightarrow \textrm{ Forme canonique Disjonctive}\\
		G_1 &=& \Pi M (1,5) \Rightarrow \textrm{ Forme Canonique Conjonctive}
	\end{eqnarray*}
	
	\chapter{Les circuits combinatoires}
	\section{Exercice 1}
	\subsection{Encodeur de priorité}
	\begin{tabular}{cccc||cc}
		$E_3$ & $E_2$ & $E_1$ & $E_0$ & $S_1$ & $S_0$\\
		\hline
		0&0&0&1&0&0\\
		0&0&1&*&0&1\\
		0&1&*&*&0&1\\
		1&*&*&*&1&1\\
	\end{tabular}\\
	\begin{eqnarray*}
	S_1&=& E_3E_2+E_3=E_3+E_2\\
	S_0 &=& E_3E_2E_1+E_3 = E_3+\overline{E_2} E_1
	\end{eqnarray*}
	\remarque{Théorème d'absorption}
	\subsection{Comparateur 4 bits}
	\subsubsection{Comparateur 1 bit}
		\begin{tabular}{cc||ccc}
			A & B & S & I & E\\
			\hline
			0&0&0&0&1\\
			0&1&0&1&0\\
			1&0&1&0&0\\
			1&1&0&0&1
		\end{tabular}
		\begin{eqnarray*}
			S&=& a\overline b\\
			I&=& \overline a b\\
			E&=& \overline a \overline b + ab = a \odot b
		\end{eqnarray*}
		\remarque{ $$a \overline \oplus b \Leftrightarrow a \odot b$$ }
	\subsubsection{Comparateur 2 bits}% a_1a_0=>A b_1b_0=>B
	$$ \underbrace{a_1a_0}_A + \underbrace{b_1b_0}_B$$
	\begin{eqnarray*}
		E &=&  (a_1 \odot b_1)(a_0 \odot b_O)\\
		S &=&  a_1 \overline{b_1} + a_0\overline{b_0}(a_1 \odot b_1)\\
		I &=& \overline{a_1}b_1 + \overline{a_0}b_0(a_1 \odot b_1)
	\end{eqnarray*}
	\subsubsection{Comparateur 4 bits -- Généralisation}
	\begin{eqnarray*}
		E &=&  (a_3 \odot b_3)(a_2 \odot b_2)(a_1 \odot b_1)(a_0 \odot b_0)\\
		S &=& a_3\overline{b_3}+a_2\overline{b_2}(a_3 \odot b_3) + a_1 \overline{b_1} (a_3 \odot b_3)(a_2 \odot b_2) + a_0
		\overline{b_0}(a_3\odot b_3)(a_2\odot b_2)(a_1 \odot b_1)\\
		I &=&  \overline{a_3}b_3 + \overline{a_2}b_2(a_3 \odot b_3) + \overline{a_1} b_1 (a_3 \odot b_3)(a_2 \odot b_3) +
		\overline{a_1}b_1 (a_3 \odot b_3)(a_2 \odot b_2) +\\&& \overline{a_0}b_0 (a_3 \odot b_3)(a_2 \odot b_2)(a_1 \odot b_1)
	\end{eqnarray*}

	\section{Exercice 2}
	\begin{tabular}{cc||cccc}
		$C_1$&$C_0$&$S_3$&$S_2$&$S_1$&$S_0$\\
		\hline
		0&0&$E_3$&$E_2$&$E_1$&$E_0$\\
		0&1&0&$E_3$&$E_2$&$E_1$\\
		1&0&0&0&$E_3$&$E_2$\\
		1&1&0&0&0&$E_3$
	\end{tabular}
	\begin{eqnarray*}
		S_3 &=&  \overline{C_1}\overline{C_0}E_3\\
		S_2 &=&  \overline C_1 \overline C_0 E_2 + \overline{C_1} C_0 E_3\\
		S_1 &=&  \overline{C_1} \overline{C_0}E_2 + C_1\overline{C_0}E_3\\
		S_0 &=& \overline{C_1} \overline{C_0}E_0 + C_1 \overline{C_0}E_1+C_1\overline{C_0}E_2 + C_1C_0E_3
	\end{eqnarray*}
\end{document}






